% !PRINTED 502
\weberpage{502}
%<<CONT
which we represent as the symbolic product $Q\pi_1$. The elements
of (3) are distinct both from one another and from the elements of
$Q$; for from $\varkappa_1\pi_1 = \varkappa_2\pi_1$ would follow $\varkappa_1 = \varkappa_2$, and from
$\varkappa = \varkappa_1\pi_1$ would follow $\pi_1 = \varkappa_1^{-1}\varkappa$, which is not possible if
$\pi_1$ does not belong to $Q$. The elements of $Q\pi_1$ form no
group, for otherwise $\varkappa_1\pi_1\varkappa_2\pi_1 = \varkappa_3\pi_1$ or $\pi_1 = \varkappa_1^{-1}\varkappa_3\varkappa_2^{-1}$
would have to hold, which would again require that $\pi_1$ be contained in $Q$.

We call the system $Q\pi_1$ a \gesperrt{subsidiary group} to $Q$
(within $P$). The subsidiary group remains unaltered in its
totality when $\pi_1$ is replaced by any element $\varkappa\pi_1$ of $Q\pi_1$,
because $Q\varkappa$ is identical with $Q$.

If we call every system $Q\pi$ generally a subsidiary group,
so that $Q$ itself is counted among them, then we can show
that two subsidiary groups are either \gesperrt{entirely identical}, or
\gesperrt{contain no common element.} For if $Q\pi_1$,
$Q\pi_2$ are two subsidiary groups, and if $\varkappa_1\pi_1 = \varkappa_2\pi_2$, in which $\varkappa_1, \varkappa_2$ are two
elements from $Q$, then, if we set $\varkappa = \varkappa_2^{-1}\varkappa_1$, it follows that
% [sic] The source prints a bare \pi here; the derivation requires \pi_1. The source wording is preserved.
$\pi_2 = \varkappa\pi$. But since $\varkappa$ belongs to $Q$, and consequently $Q\varkappa = Q$,
it follows that
\[
Q\pi_2 = Q\varkappa\pi_1 = Q\pi_1.
\]
If the group $P$ is not yet completely exhausted by (2) and (3), we can form a further subsidiary group $Q\pi_2$,
and can continue in the formation of the subsidiary groups until $P$ is decomposed into nothing but subsidiary groups. Since each of these
% [sic] The source gives the number of elements in each subsidiary group as \nu rather than the previously defined group degree q; preserved source-faithfully.
subsidiary groups contains $\nu$ elements, there follows the important fundamental theorem of \gesperrt{Cauchy}:

1. \gesperrt{The degree of a divisor of a group $P$ is
a divisor of the degree of $P$.}

As a special case there is contained in this the fact that the \gesperrt{degree
of an element of the group} (\S.\,153) \gesperrt{is always a divisor
of the degree of the group}.

The decomposition of $P$ into its subsidiary groups we can
(following \gesperrt{Galois}) express very fittingly by the symbolic equation
\weq{4}{P = Q + Q\pi_1 + Q\pi_2 + \cdots + Q\pi_{j-1},}
if
\weq{5}{p = jq}
holds. The quotient $j$ we call the \gesperrt{index of the divisor}
$Q$. This number expresses no peculiarity of the group $Q$
in itself, but only a relation of $Q$ to $P$.
